\documentclass[10pt]{article}
\usepackage[T1]{fontenc}
\usepackage[utf8]{inputenc}
\usepackage[french]{babel}
\usepackage{amsmath,amssymb,mathrsfs,mathpazo}
\usepackage{geometry}
\geometry{a4paper,margin=2cm}
\usepackage{graphicx,array,tabularx,tabularray,colortbl,multirow,multicol,booktabs,enumitem,pifont,ulem,fancybox,diagbox,cancel}
\usepackage{pstricks,pstricks-add,pst-plot,pst-tree,pst-node,pst-text,pst-func,pst-math,pst-bspline,pst-3dplot,pst-eucl,pst-solides3d}
\let\pscalculate\undefined
\usepackage{pst-bezier}
\usepackage{tikz,pgfplots,tkz-tab}
\usetikzlibrary{arrows,arrows.meta,patterns,calc,positioning,trees,decorations.markings,shapes.geometric,shapes.misc,shapes.arrows,intersections,decorations.pathmorphing,decorations.pathreplacing,backgrounds,shadows,babel}
\pgfplotsset{compat=1.18}
\usepgfplotslibrary{fillbetween}
\usepackage[np]{numprint}
\usepackage{listings,makecell,pas-tableur,twoopt}
\let\barre\undefined
\usepackage[pstricks]{tablvar}
\definecolor{darkgreen}{RGB}{0,100,0}
\providecommand{\red}{\color{red}}
\providecommand{\blue}{\color{blue}}
\providecommand{\green}{\color{green}}
\providecommand{\black}{\color{black}}
\providecommand{\white}{\color{white}}
\providecommand{\gray}{\color{gray}}
\providecommand{\cyan}{\color{cyan}}
\providecommand{\magenta}{\color{magenta}}
\providecommand{\yellow}{\color{yellow}}
\providecommand{\violet}{\color{violet}}
\providecommand{\orange}{\color{orange}}
\providecommand{\euro}{\texteuro}
\providecommand{\EUR}{\texteuro}
\providecommand{\ieme}{\textsuperscript{e}}
\providecommand{\degres}{\ensuremath{^\circ}}
\providecommand{\degree}{\ensuremath{^\circ}}
\providecommand{\vv}[1]{\overrightarrow{#1}}
\providecommand{\Coord}[1]{#1}
\providecommand{\decosix}{}
\providecommand{\starredbullet}{\ensuremath{\bullet}}
\pagestyle{empty}
\setlength{\parindent}{0pt}
\usepackage{adjustbox}\DeclareUnicodeCharacter{2236}{:}\DeclareUnicodeCharacter{2190}{\ensuremath{\leftarrow}}
\providecommand{\R}{\mathbb{R}}
\providecommand{\N}{\mathbb{N}}
\providecommand{\Z}{\mathbb{Z}}
\providecommand{\Q}{\mathbb{Q}}
\providecommand{\D}{\mathbb{D}}
\providecommand{\C}{\mathbb{C}}
\providecommand{\e}{{\mathrm{e}}}
\providecommand{\vect}[1]{\overrightarrow{#1}}
\providecommand{\vectt}[1]{\overrightarrow{\mathrm{#1}}}
\providecommand{\vvt}[1]{\overrightarrow{\mathrm{#1}}}
\providecommand{\np}[2][]{{#2~#1}}
\providecommand{\vv}[1]{{\overrightarrow{#1}}}
\renewcommand{\d}{{\mathrm{d}}}
\providecommand{\degres}{{^\circ}}
\providecommand{\degre}{{^\circ}}
\providecommand{\blue}{\color{blue}}
\providecommand{\red}{\color{red}}
\providecommand{\enskip}{\quad}
\providecommand{\llbracket}{\lbrack\!\lbrack}
\providecommand{\rrbracket}{\rbrack\!\rbrack}
\providecommand{\euro}{{\text{€}}}
\providecommand{\no}{{\text{n°}}}
\providecommand{\interval}[2]{{[#1;#2]}}
\providecommand{\intervFO}[2]{{[#1;#2[}}
\providecommand{\dotfill}{\ldots}

\providecommand{\ds}{\displaystyle}
\providecommand{\textding}[1]{#1}
\providecommand{\Oij}{\(\mathrm{O};\vec{\imath},\vec{\jmath}\)}
\providecommand{\Oijk}{\(\mathrm{O};\vec{\imath},\vec{\jmath},\vec{k}\)}
\providecommand{\Ouv}{\(\mathrm{O};\vec{u},\vec{v}\)}
\providecommand{\intervFF}[2]{\ensuremath{[#1;#2]}}
\renewcommand{\euro}{\ifmmode\text{\texteuro}\else\texteuro\fi}
\DeclareUnicodeCharacter{00B0}{\ifmmode^{\circ}\else\textdegree\fi}
\DeclareUnicodeCharacter{20AC}{\ifmmode\text{\texteuro}\else\texteuro\fi}
% Énoncés : Éducation nationale. Transcription et corrigés : APMEP, Denis Vergès et contributeurs. Découpage et classement : CoopMaths, Sébastien Lozano, Éric Elter et contributeurs. Adaptation de la présentation : Maths Explorer.
% CC BY-SA : https://coopmaths.fr/www/about/
% https://coopmaths.fr/alea/static/bac/2026/tex/bac_2026_06_sujet1_metropole_2_cor.tex
\begin{document}
\noindent\textit{APMEP / CoopMaths. CC BY-SA. Présentation adaptée par Maths Explorer.}\par\medskip
\begin{enumerate}
		\item \begin{enumerate}
			\item \textbf{Affirmation 1 : Vraie.}

			Par lecture de l'équation cartésienne du plan on peut affirmer que $\vect{n}\begin{pmatrix}-1\\1\\-5\end{pmatrix}$ est un vecteur normal au plan $(\mathcal{P})$.

			La droite (AB) est dirigée par le vecteur $\vectt{AB}$, de coordonnées : \quad $\begin{pmatrix} x_{\mathrm{B}} - x_{\mathrm{A}} \\ y_{\mathrm{B}} - y_{\mathrm{A}} \\ z_{\mathrm{B}} - z_{\mathrm{A}} \end{pmatrix} $

			soit :\quad $\vectt{AB} \begin{pmatrix} 2 - 3\\1-0\\3-2 \end{pmatrix} $ \quad et donc :\quad $\vectt{AB} \begin{pmatrix} -1\\1\\-5 \end{pmatrix} $.

			Comme $\vectt{AB} = \vect{n}$, le vecteur directeur de la droite (AB) est donc normal au plan $(\mathcal{P})$.

			Par définition, la droite (AB) est donc orthogonale au plan $\mathcal{P}$.

			Le milieu du segment [AB] est le point I, et ses coordonnées vérifient :

			$\mathrm{I}\left(\dfrac{x_{\mathrm{A}} + x_{\mathrm{B}}}{2} ; \dfrac{y_{\mathrm{A}} + y_{\mathrm{B}}}{2} ; \dfrac{z_{\mathrm{A}} + z_{\mathrm{B}}}{2}\right) $\quad soit \quad $\mathrm{I} \left(\dfrac{3+2}{2} ; \dfrac{0+1}{2} ; \dfrac{2-3}{2}\right) $\quad et finalement \quad $\mathrm{I}(2,5 ; 0,5 ; -0,5)$.

			Vérifions si les coordonnées de I vérifient l'équation du plan $(\mathcal{P})$ :

			$-x_{\mathrm{I}}+y_{\mathrm{I}}-5 z_{\mathrm{I}}-0,5= -2,5 +  0,5 -5 \times (-0,5) - 0,5 = -2,5 +0,5 + 2,5 - 0,5 = 0$

			I est donc un point du plan $(\mathcal{P})$.

			Le plan $(\mathcal{P})$ est bien orthogonal à la droite (AB), et passe par le milieu du segment [AB]

			(autrement dit, c'est le plan médiateur du segment [AB]).

			\item \textbf{Affirmation 2 : Fausse.}

			La droite (AB) est dirigée par le vecteur $\vectt{AB} \begin{pmatrix} -1\\1\\-5 \end{pmatrix} $.

			Par lecture de sa représentation paramétrique, on peut dire que $(d)$ est dirigée par $\vect{u}\begin{pmatrix} 1\\-1\\-2 \end{pmatrix}$.

			Les deux vecteurs ont des coordonnées clairement non proportionnelles, donc ils ne sont pas colinéaires.

			Les droites (AB) et $(d)$ ne sont donc pas parallèles (ni strictement, ni confondues).

			Elles peuvent donc être sécantes, ou non coplanaires.

			Pour trancher, cherchons à établir si elles ont un point commun.

			Une représentation paramétrique de (AB) est :\quad $\begin{cases}
				x = x_{\mathrm{A}} + x_{\vectt{AB}}\times s\\
				y = y_{\mathrm{A}} + y_{\vectt{AB}}\times s\\
				z = z_{\mathrm{A}} + z_{\vectt{AB}}\times s\\
			\end{cases}\quad \text{ où } s \in \R$

			Soit, en remplaçant par les valeurs connues : \quad $\begin{cases}
				x = 3 - s\\ y = s \\ z = 2 - 5s
			\end{cases}\quad \text{ où } s \in \R$

			Le point de paramètre $t$ sur $(d)$ est confondu avec le point de paramètre $s$ sur (AB) si et seulement si $(s ; t)$ est solution du système :

			$\begin{aligned}
				\begin{cases}
					3-s = t\\
					s = -1,5 - t\\
					2 -5s = 2 - 2t
				\end{cases}
				&\iff \begin{cases}
					3-s = t\\
					s = -1,5 - (3 - s)\\
					2 -5s = 2 - 2(3 - s)
				\end{cases}\\
				&\iff \begin{cases}
					3-s = t\\
					s = -1,5 - 3 + s\\
					2 -5s = 2 - 6 + 2s
				\end{cases}\\
			\end{aligned} $

			Les équations du système sont incompatibles (la deuxième équation est équivaut à $0 = -4,5$).

			Le système n'a pas de solution.

			Il n'y a donc aucun point commun entre les droites (AB) et $(d)$, qui sont donc non coplanaires, et pas sécantes.

			\item \textbf{Affirmation 3 : Vraie}

			On a :\quad $\vectt{CA} \begin{pmatrix} 1,5\\3\\3 \end{pmatrix} $\quad et \quad $\vectt{CB} \begin{pmatrix} 0,5\\4\\-2 \end{pmatrix} $.

			Comme ce sont des coordonnées dans un repère orthonormé, on utilise la formule analytique du produit scalaire :\quad $\vectt{CA}\cdot\vectt{CB} = 1,5 \times 0,5 + 3 \times 4 + 3 \times (-2) = 0,75 + 12 - 6 = 6,75$.

			Par ailleurs, on peut calculer les longueurs CA et CB :

			$\mathrm{CA} = \sqrt{1,5^2 + 3^2 + 3^2} = \sqrt{20,25} = 4,5$.\quad $\mathrm{CB} = \sqrt{0,5^2 + 4^2 + (-2)^2} = \sqrt{20,25} = 4,5$.

			On remarque donc (même si c'est hors sujet) que le triangle ACB est isocèle en C.

			En utilisant la définition trigonométrique du produit scalaire, on a :

			$\vectt{CA}\cdot\vectt{CB} = \mathrm{CA} \times \mathrm{CB} \times \cos\left(\widehat{\mathrm{ACB}}\right)$.

			En remplaçant les valeurs connues, on a :\quad $6,75 = 4,5 \times 4,5 \times \cos\left(\widehat{\mathrm{ACB}}\right)$,

			ce qui équivaut à \quad $\cos\left(\widehat{\mathrm{ACB}}\right) = \dfrac{6,75}{4,5^2} = \dfrac{1}{3}$.

			On en déduit :\quad $\widehat{\mathrm{ACB}} = \arccos\left(\dfrac{1}{3}\right) \approx 70,53$\,\textdegree.

			On a donc bien $\widehat{\mathrm{ACB}} \approx 70,5$\,\textdegree, à $10^{-1}$ près.
		\end{enumerate}

		\item \textbf{Affirmation 4 : Vraie.}

		\begin{itemize}[label=\textbullet~]
			\item Titouan doit, pour ouvrir sa porte, trouver les quatre symboles différents de son code, et les saisir dans un ordre quelconque.

			Il s'agit donc d'une combinaison (pas de répétition des éléments, l'ordre ne compte pas) de quatre éléments, dans un ensemble de cardinal 8 (huit symboles sur le digicode).

			$\ds \binom{8}{4} = 70$.

			Titouan a donc 70 codes différents possibles, chacun ayant donc une probabilité de $\dfrac{1}{70}$, puisque tous les choix des codes sont équiprobables.
			\item Clotilde, elle, doit saisir trois symboles différents, mais en tenant compte de l'ordre.

			Il s'agit donc d'un arrangement (pas de répétition, l'ordre compte) de trois éléments dans un ensemble de cardinal 8.

			$ \dfrac{8!}{(8-3)!} = \dfrac{8!}{5!} = 8\times 7 \times 6 = 336$.

			Clotilde a donc 336 codes différents possibles, tous ayant une probabilité de $\dfrac{1}{336}$.
		\end{itemize}

		Comme \quad $\dfrac{1}{70} > \dfrac{1}{336}$\quad la probabilité que Titouan ouvre sa porte est bien supérieur à la probabilité que Clotilde ouvre la sienne.
	\end{enumerate}

\end{document}
